%*****************************************
\chapter{Input e output}
%*****************************************

Da qua in avanti potrebbe capitare con crescente frequenza di incontrare nuovamente esercizi già visti, più o meno simili. Non è un errore o una svista: molti esercizi sono effettivamente riciclati perché possono essere ripensati in modo via via più completo utilizzando i nuovi strumenti appresi.

Alcuni consigli utili per capire come identificare e gestire i dati in un problema:

\begin{aenumerate}
	\item quando si identificano i dati nel problema, distinguere subito quali sono utili alla sua risoluzione e quali no poiché alcuni potrebbero essere superflui o distrattori;
	\item alcuni dati sono già noti nella descrizione del problema, altri invece vanno chiesti in input, ed è sempre bene distinguere chiaramente queste due categorie di dati;
	\item i dati letti in input vengono memorizzati nelle variabili e possono poi essere trattati come tutti gli altri dati.
\end{aenumerate}

\section{Qualche esercizio di riscaldamento}

\begin{exercise}
	Realizzare un algoritmo che legga in input la misura del lato di un quadrato e ne calcoli il perimetro.
\end{exercise}

\begin{exercise}
	Realizzare un algoritmo che legga in input la misura del raggio di una circonferenza e ne calcoli l'area.
\end{exercise}

\begin{exercise}
	Realizzare un algoritmo che calcoli il resto di un pagamento sapendo che il pacco di Nutella biscuits costa 2,68 € e leggendo in input la somma usata per pagare, alla fine stampare in output il resto.
\end{exercise}

\begin{exercise}
	Realizzare un algoritmo che calcoli i minuti necessari a percorrere una distanza in chilometri ricevuta in input sapendo che si viaggia ad una velocità media di 90 km/h, alla fine stampare in output il tempo calcolato.
\end{exercise}

\begin{exercise}
	Realizzare un algoritmo che calcoli e stampi in output quanti atomi di idrogeno ci sono in totale in un certo numero di molecole di acqua letto in input sapendo che in una molecola di acqua sono sempre presenti due atomi di idrogeno e uno di ossigeno.
\end{exercise}

\begin{exercise}
	Realizzare un algoritmo che calcoli e stampi in output il numero di uova prodotte in un pollaio dopo un certo numero di giorni letto in input sapendo che: nel pollaio ci sono 17 galline e che il pollaio produce mediamente 15 al giorno.
\end{exercise}

\begin{exercise}
	Realizzare un algoritmo che stampi in output l'età di una persona dopo che è stato inserito in input il suo anno di nascita.
\end{exercise}

\begin{exercise}
	Realizzare un algoritmo che calcoli i consumi idrici di una famiglia sapendo che un italiano in media consuma 220 litri al giorno; si legga in input il numero di membri della famiglia e quindi si stampi in output il consumo della famiglia dopo una settimana.
\end{exercise}

\begin{exercise}
	Realizzare un algoritmo che calcoli quanti sacchi di crocchette comprare per i propri gatti sapendo che un gatto mediamente mangia 50 grammi di crocchette e i sacchi sono da 0,7 kg, chiedere in input il numero di giorni e stampare in output quanti sacchi di crocchette si dovrebbero comprare.
\end{exercise}

\section{Più input e più output}

\begin{exercise}
	Realizzare un algoritmo che calcoli e stampi in output le dimensioni in metri quadri di una parete chiedendo in input altezza e larghezza.
\end{exercise}

\begin{exercise}
	Realizzare un algoritmo che calcoli e stampi in output le dimensioni in metri quadri di una parete chiedendo in input altezza e larghezza; sapendo che con un litro di vernice si possono tinteggiare 4 mq, calcolare e stampare in output anche i litri di vernice necessari per tinteggiare l'intera parete.
\end{exercise}

\begin{exercise}
	Realizzare un algoritmo per il calcolo dei consumi idrici di una famiglia leggendo in input il numero di membri che la compongono. In media un italiano consuma 220 l di acqua al giorno, calcolare i consumi dopo una settimana esprimendo il risultato sia in litri che in metri cubi.
\end{exercise}

\begin{exercise}
	Realizzare un algoritmo che, letta in input una distanza percorsa espressa in chilometri e il tempo impiegato espresso in minuti, calcoli e stampi in output la velocità espressa in km/h.
\end{exercise}

\begin{exercise}
	Realizzare un algoritmo che legga in input il numero di brioche e il numero di caffè presi a colazione, quindi calcoli e stampi il costo totale sapendo che una brioche costa 1,40 € e un caffè 1,10 €.
\end{exercise}

\begin{exercise}
	Realizzare un algoritmo che legga in input quante galline e quanti galli ci sono in un pollaio, quindi calcoli la produzione mensile di uova e il numero di confezioni che servono per vendere le uova. In media due galline producono 7 uova a settimana, le confezioni che si usano per le uova sono da 6 e per semplicità si asumma che un mese è fatto da 4 settimane.
\end{exercise}

\begin{exercise}
	Realizzare una algoritmo che calcoli quanti ingredienti comprare per una pizzeria chiedendo in input il numero di pizze da produrre. Il pizzaiolo stima che per una pizza servano 150 g di farina, 80 g di acqua, 9 g di lievito e 10 g di olio. Stampare i quantitativi in kg di ogni ingrediente da acquistare.
\end{exercise}

\begin{exercise}
	Realizzare un algoritmo che legga in input le due dimensioni di un rettangolo e ne calcoli perimetro e area.
\end{exercise}

\begin{exercise}
	Realizzare un algoritmo che legga in input l'area di un quadrato e ne calcoli lato e perimetro.
\end{exercise}

\begin{exercise}
	Realizzare un algoritmo che calcoli e stampi la superficie e il volume di un parallelepipedo dopo aver letto in input le misure delle sue tre dimensioni: larghezza, altezza e profondità.
\end{exercise}

\begin{exercise}
	Realizzare un algoritmo che calcoli e stampi la superficie e il volume di un cilindro dopo aver letto in input i dati necessari.
\end{exercise}

